\section{Sequencer}\label{sequencer}

The sequencer plays the Pattern of each track's played Scene. The basic
idea is familiar from classic step machines, with considerably more
per-step control.

\subsection{Tracks}\label{tracks}

There are \textbf{7 tracks}: one per voice, plus an extra track that
plays the open hi-hat of voice 6. Each track has its own length, step
scale, shuffle and run mode, so tracks can drift against one another and
realign.

\subsubsection{Select, mute and unmute}\label{select-mute-and-unmute}

\begin{itemize}
\tightlist
\item
  The active track is chosen with the 7 \textbf{voice buttons}. Its LED
  is lit.
\item
  \textbf{SHIFT + voice button} mutes or unmutes the track.
\item
  In Performance mode the roles are inverted: a short tap on a voice
  button toggles mute without SHIFT, and holding it allows other
  gestures (per-track Scene playback, overlay triggers).
\item
  A muted track does not play and does not play its step automation.
  Effect tails keep ringing.
\end{itemize}

\subsection{Steps}\label{steps}

Each track has up to \textbf{128 steps}, shown as eight bars of 16
steps. The 16 sequencer buttons show the bar currently being played or
viewed. Press a button to set or clear a step.

\subsubsection{Bars and sub steps}\label{bars-and-sub-steps}

The eight SELECT LEDs show the bar of the active track that is playing;
the LED of a new bar pulses once. Sub steps for fine placement are
edited in Step mode.

\figplaceholder{substep_diagram.png}{Steps, bars and sub steps on the sequencer buttons.}{Row of 16 step buttons with the LEDs of the current bar, and below it the eight SELECT LEDs showing which of the eight bars is shown. Second panel: one step expanded into its sub-step divisions.}

\verify{The earlier LXR had 16 main steps with 8 sub steps each. The new Pattern model is 7 tracks of 128 steps, shown as 16 steps per bar. Confirm how sub steps are presented in Step mode and update this section and the figure.}

\subsection{Per-track settings (the STEP
page)}\label{per-track-settings-the-step-page}

Open Step mode and press the encoder to the track settings page. Page 1
has \texttt{len\ scl\ shf\ run}; page 2 has \texttt{ch\ nte\ lop\ rmp}.

\begin{paramtable}
len & Length      & Number of steps in the track (1 to 128). Each track may have a different length. \\
scl & Step scale  & How fast the track steps. A continuous scale from 1/64 (0) through 1/16 (76, the default) to 1/2 (127), with named stops: \texttt{/64 32t /32 16t d32 /16 8Tr d16 /8 4Tr d/8 /4 d/4 /2}. Other positions show the raw number. \\
shf & Shuffle     & Delays every second step on the 1/16 grid. Independent of scale. \\
run & Run mode    & \texttt{fwd} forward, \texttt{rev} reverse, \texttt{pip} ping-pong (end steps play twice), \texttt{rnd} random, \texttt{onc} once from the master clock then stop, \texttt{1fr} once from step 1 then stop. \\
ch  & MIDI channel & \texttt{off}, \texttt{glb} (follow the global channel), or 1 to 16. \\
nte & MIDI note   & \texttt{off}, \texttt{on} (no root note), or the root note of this track's receive range. \\
lop & Loopable    & Whether the PERF looper may loop this track. \\
rmp & Automation ramp & On: step automation values glide from step to step instead of jumping. \\
\end{paramtable}

Length, step scale and shuffle are also Morph targets (hold SHIFT on the
page to edit the Morph endpoint) and can be automated per step. Changing
the scale mid-playback never causes drift.

\verify{Confirm the abbreviation of the Loopable cell and that SHIFT shows Morph endpoints for len scl shf only.}

\subsection{Step parameters}\label{step-parameters}

Hold a step and edit with the knobs (or select the step in Step mode) to
reach its \textbf{specials}. Every special is optional; a step with no
specials plays at the default velocity and pitch.

\begin{paramtable}
vel & Velocity   & Note-on velocity, 0 to 127. Velocity 0 is an \emph{automation-only} step: envelopes are not retriggered but automation still applies. \\
nte & Note       & Pitch offset in semitones relative to the default pitch. \\
cnd & Condition  & Whether the step plays; see below. \\
rol & Roll       & Retrigger rate: \texttt{off 1/8 8tr 16 16t 32 32t 64 64t 128}. \\
tim & Timing     & Micro-timing: delays the step by a percentage of the step length. \\
rtm & Rhythm     & Velocity pattern for the repeats of a roll. \\
fil & Fill       & Percentage of the step length that the roll lasts (100 is the full step). \\
\end{paramtable}

\verify{On-screen tokens for condition, timing and fill are taken from the data model; only rol, rtm and the roll labels are confirmed in the menu code.}

\subsubsection{Conditions}\label{conditions}

A step with a condition plays only when the condition is true. A failed
condition suppresses the trigger \textbf{and} the step automation. The
default is always-play.

\begin{itemize}
\tightlist
\item
  \textbf{Probability:} play the step with a fixed chance (ten levels).
\item
  \textbf{Repeat X:Y:} play on the X-th of every Y passes through the
  track (ten combinations).
\item
  \textbf{Fill / not fill:} play only while a fill is held (hold the
  Performance fill gesture longer than 300 ms), or only while no fill is
  held.
\item
  \textbf{Loop / not loop:} play only while the track is being looped by
  the PERF looper, or not.
\item
  \textbf{Last / not last:} depends on whether the previous conditioned
  step fired.
\item
  \textbf{Mute / not mute of track 1 to 7:} play only when the named
  track is muted (or not). A step that tests its own track's mute state
  is never silenced by that mute.
\end{itemize}

\tip{Conditions are evaluated once per step. A roll on the step uses the same decision for all
its repeats.}

\subsubsection{Rolls}\label{rolls}

A step with a roll special retriggers at the chosen rate for as long as
its \textbf{window} lasts, which is the step length multiplied by the
fill percentage. The rate is an absolute note value; the track scale
decides only how long the window is, never the spacing of the hits.

\begin{itemize}
\tightlist
\item
  Two adjacent roll steps with the \textbf{same rate and rhythm} and a
  fill of 100 are \emph{carried}: the roll continues across the boundary
  without an extra hit and without losing its rhythm.
\item
  A step whose fill is below 100 ends its roll with a short gate duck,
  and the next step starts a fresh roll.
\item
  A plain hit ends a running roll; a different roll replaces it. An
  empty or automation-only step does not interrupt it.
\item
  The \textbf{rhythm} (\texttt{rtm}) shapes the velocities of the
  repeats: \texttt{off}, \texttt{1-2} to \texttt{4-5} (every n-th repeat
  is halved) or \texttt{1+2} to \texttt{4+5} (every n-th repeat is
  doubled).
\end{itemize}

\subsection{Automation}\label{automation}

Each step can carry \textbf{parameter automation}: any Instrument,
Effect or Scene parameter that is marked automatable. Automation affects
only that step and is ignored if the step does not play.

\subsubsection{Priority rules}\label{priority-rules}

\begin{itemize}
\tightlist
\item
  \textbf{Instrument parameters} (rule R12): a step's value holds until
  the voice's next \emph{first} trigger of a step. Roll repeats after
  the first hit never reset it. A \textbf{menu edit}, including a MIDI
  write, overrides the automation immediately and keeps it suppressed
  until the voice's next trigger, a transport change or a Scene change.
\item
  \textbf{Effect and Scene parameters} (rule R13): the last automation
  sets the value and it stays until the next automation; an
  effect-sequence lock beats track automation for its own step; a menu
  edit beats both until the next automation.
\item
  With the track's \textbf{\texttt{rmp}} setting on, numeric parameters
  ramp linearly between consecutive played steps that both hold them.
  Ramps end at the next played step and at every position break, and
  never cross a loop wrap.
\item
  The LFO and Morph modulate \textbf{around} whatever value automation
  sets.
\end{itemize}

\subsubsection{Adding automation}\label{adding-automation}

\textbf{Hold a step and turn a knob} (or encode a value on the Effect
page): the value is locked on that step. The step's LED flashes while it
is held. Hold several steps to write the same value to all of them.

\textbf{Edit in Step mode:} the step's automation page lists its
entries. Choose a destination category and parameter and set a value.
Each step holds a number of entries, each with a unique target.

\textbf{Record live:} see \emph{Live Recording} in the Operation Modes
chapter.

\subsubsection{Underlines}\label{underlines}

A parameter name is underlined when any step automation addresses it.
The search runs in the background, so underlines can appear a moment
after you enter a page.

\subsection{Per-track Scene playback}\label{per-track-scene-playback}

In Performance mode, hold a voice button and press a SEQ button to make
that track play a different Scene's Pattern, instrument parameters, MIDI
routing, fader and effect send. A double-click on a SEQ button toggles
the override. Tracks 6 and 7 move together. Recording is disabled while
any track plays a different Scene.

Pressing the pattern button of the currently viewed Scene again realigns
every track to the master clock, which is useful after tracks with
different scales have drifted.

\subsection{Copy and Clear}\label{copy-and-clear}

\draftnote{The Copy/Clear specification is 85 KB; this section is a summary. Verify every gesture
on the device.}

Copy and clear work at several levels: single steps and ranges, bars,
tracks, automation, whole Patterns, and Scenes and their children (Kit,
Effect, Morph). The Instrument and Effect level operations follow the
voice edit mask.

\subsubsection{Copy}\label{copy}

\begin{enumerate}
\def\labelenumi{\arabic{enumi}.}
\tightlist
\item
  Hold \textbf{Copy}; its LED flashes.
\item
  Press the source (a voice button for a track, a SEQ button for a step,
  a SELECT button for a bar or sub step).
\item
  Press the destination. The data is pasted at once and overwrites the
  destination.
\item
  Release Copy to abort.
\end{enumerate}

\subsubsection{Clear}\label{clear}

\begin{enumerate}
\def\labelenumi{\arabic{enumi}.}
\tightlist
\item
  Press \textbf{SHIFT + Copy/Clear}.
\item
  The display offers what to clear (for example track, pattern,
  automation, Effect, Scene, Morph). Use the encoder to select.
\item
  Press Copy/Clear again to confirm, or release SHIFT to abort.
\end{enumerate}

While a copy or clear is running, background saving is suspended until
it finishes.

\verify{Exact choice list shown by the clear menu in each mode.}
